\documentclass[10pt,letterpaper]{article}
\usepackage[margin=0.85in]{geometry}
\usepackage[T1]{fontenc}
\usepackage{lmodern}
\usepackage{amsmath,amssymb,booktabs,array}
\usepackage[hidelinks]{hyperref}
\usepackage{microtype}
\setlength{\parskip}{3pt}
\setlength{\parindent}{1em}
\newcommand{\claim}[1]{\textsuperscript{\scriptsize[#1]}}
\newcommand{\todo}[1]{\textit{Citation TODO: #1.}}
\title{Joint Interpretation of Time-Varying Factual Records\\
\large Exploratory manuscript skeleton for an ACL 2027 research target}
\author{Working draft; authorship to be supplied}
\date{2 October 2026 \quad | \quad W01 / v0.7}

\begin{document}
\maketitle
\begin{abstract}
Public records can announce future changes, describe an ongoing state, repeat an
earlier assertion, or correct what was previously reported. We study a proposed
task in which a system selects interpretations of source mentions together with
their temporal and correction relations, then answers questions conditional on
the resulting record. This draft specifies the implemented small-instance
reference semantics, shared inference objective, and comparison protocol. It
also documents two development observations: agreement among the implemented
decoders on one model-scored natural example, and correction suppression under
an active-support objective ablation with existing authored scores. Neither
observation establishes a natural-data accuracy gain. Independent evaluation,
strong contextual baselines, candidate-coverage assessment, and a broader
related-work audit remain necessary before a paper contribution can be stated.
\end{abstract}

\noindent\textbf{Draft status.} This is a methods and data-design skeleton, not a
submission-ready paper. Numbered claim markers refer to the accompanying
\texttt{claim\_ledger\_v07.json}. Results here are frozen to the stated v0.6
artifacts; ongoing v0.7 backend and scoring work is not reported as complete.
The article layout is portable LaTeX, not an asserted ACL 2027 format.

\section{Research question and information boundary}
\label{sec:task}
The working question is whether revisiting uncertain factual interpretations
together with their update relations improves correction-dependent record
reconstruction beyond strong contextual extraction and conditional relation
inference. An interpretation can
change a subject, relation, scope, value, episode assignment, temporal boundary,
or correction target. These are distinct error categories. Their frequency and
their susceptibility to joint inference are empirical questions, not assumed
properties of every factual history.\claim{C01}

Let a source record $d$ have text, a version identifier, and a recorded
availability day $a_d$. At information cutoff $\tau$, the eligible source set is
$D_\tau=\{d:a_d\leq\tau\}$. Every declared context dependency of a candidate
must also belong to this set. A question's event day $t$ is separate from
$\tau$: the system may ask what a record available in September announced for
October, without treating that announcement as evidence of the event's eventual
occurrence. A claim's report date is another field and does not automatically
mark an event onset or a state observation.\claim{C02}

The implementation checks declared provenance, dates, and exact source spans.
It cannot authenticate an assigned historical availability date or detect an
undeclared context dependency. A current page and a retrieved text hash establish
the captured bytes, not when those bytes first became public. Development
eligibility conventions must therefore remain explicit, and current-record
experiments must not be described as historically faithful prefix evaluation.
The original and revised information prefixes are constructed separately.

The prospective primary study instead defines a controlled stream of frozen
document versions. At delivery batch $b$, its prefix is
$D^{\mathrm{stream}}_b=\{d:\operatorname{batch}(d)\leq b\}$. Delivery index,
retrieval time, reported publication time, and event time remain separate.
This task asks what the supplied prefix justifies, not what all public evidence
justified at a historical instant. A separately certified historical slice
would require evidence for the original versions and their availability bounds.
No histories have yet been completed under this new stream protocol.\claim{C12}

\section{Candidate representation and conditional semantics}
\label{sec:representation}
For each anchored source mention $m$, the candidate set $R_m$ contains resolved
readings and at least one unresolved reading. A resolved reading supplies a
factual key $k=(\text{subject},\text{relation},\text{scope})$, value, polarity,
modality, evidence spans, and bounded temporal information. Scope must carry
distinctions needed for a single-valued relation; a descriptive role qualifier
does not silently create a separate key. Unresolved readings do not invent a
value or exact date. An assignment selects exactly one $r_m\in R_m$ and one
outgoing typed link $\ell_m$ or null per mention.\claim{C03}

Candidates are hypotheses, not automatically accepted source interpretations.
Syntactic integrity and exact span alignment do not prove that a quote entails
a candidate. In particular, a reading that is plausible in a local passage may
be rejected by other available context. Such cases can test inference even if
the full document leaves no residual ambiguity. A separate audit that requires
two defensible full-context readings measures a narrower phenomenon and cannot
establish candidate recall or the absence of useful joint reasoning.\claim{C12}

\subsection{Episodes and typed temporal links}
\label{sec:temporal}
An active episode has integer-day endpoints $s,e$ and denotes the half-open
interval $[s,e)$, with $s+1\leq e$. Supplied endpoint ranges are inclusive bounds
on the latent endpoints. An observation on day $t_o$ adds
$s\leq t_o$ and $t_o+1\leq e$. Missing bounds remain unknown; they are not a
persistence rule. Combining observations into one episode does assume
continuity between them.\claim{C04}

\textsc{Restates} asserts equality of the two episode endpoints and requires
equal keys, values, polarity, and modality. Equal values alone do not imply the
same episode: a return to value A after an intervening B can be a distinct A
episode. \textsc{Changes} points from a later episode to an earlier episode,
requires equal keys, distinct positive values, and equal modality, and imposes
$e_{\mathrm{earlier}}\leq s_{\mathrm{later}}$. It permits a gap and does not
assert immediate succession. Null adds no temporal relation. Unlinked claims
can still overlap or conflict; the solver does not impose universal same-key
nonoverlap.\claim{C04}

Each inequality $x_j-x_i\leq c$ becomes an edge $i\to j$ of weight $c$, with a
fixed zero variable for calendar bounds. Shortest-path closure detects an
infeasible state by a negative cycle and projects tight endpoint bounds.
For an active episode with projected bounds
$s\in[s_l,s_u]$ and $e\in[e_l,e_u]$, conditional membership is
\begin{align}
\operatorname{may}(t)&\iff s_l\leq t<e_u, &
\operatorname{must}(t)&\iff s_u\leq t<e_l. \label{eq:membership}
\end{align}
These statements are conditional on the selected readings and links. They do
not quantify over alternative decoder assignments. The temporal-network solver
and exhaustive inference are reference techniques; no novelty is claimed for
them. This feasibility layer instantiates the established temporal constraint
network family~\cite{dechter1991temporal}.

\subsection{Correction and essential support}
\label{sec:correction}
\textsc{Corrects} replaces support for a particular earlier atomic assertion;
it does not end that assertion's real-world interval on the correction date.
A selected correction requires an exact reference span, compatible factual
keys, recorded nondecreasing source availability, and equal annotated asserting
authority. Hosting domains do not establish authority. These conditions check
the representation of an attributed correction, not independently verified
entailment or speaker identity.\claim{C05}

The selected correction graph must be acyclic. Its targets are withdrawn before
temporal closure. In a chain where B corrects A and C corrects B, both A and B
stay withdrawn: withdrawal of B's replacement content does not reverse B's
correction notice. Declared essential-support dependencies then propagate
inactivity to a fixed point if a prerequisite is absent or inactive. Dependencies
are conjunctive and are not inferred merely from a \textsc{Restates} relation.
Inactive assertions and their incident temporal links supply no temporal bounds
or answer support. Correction links supply no event-date inequality.

The implemented scope is whole-assertion replacement. Bare retraction,
partial-field revision, retraction of a correction notice, and arbitrary
belief revision require a richer representation. Independently selected sibling
corrections can remain conflicting assertions; the implementation does not
manufacture a winner from chronology alone.\claim{C05}

Maintaining reasons for beliefs and revising dependent support has a classical
truth-maintenance precedent~\cite{doyle1979truth}. Our restricted projection
specifies one correction policy; it neither invents dependency withdrawal nor
implements general truth maintenance or arbitrary belief revision.

\subsection{Answers and their interpretation}
\label{sec:answers}
Queries distinguish announced schedules from reported actual state. Announced
future claims cannot prove actual realization simply because the scheduled day
has passed. Actual-state queries preserve each claim's original report horizon,
with source availability as a documented fallback when no report date exists.
Conditional or uncertain claims can supply possibilities but not certainty.
An affirmative answer requires one certainly supported value without a surviving
possible alternative or conflicting negative support. Otherwise the memory
returns indeterminate or abstain, according to whether affirmative possibilities
remain. Negative-only evidence cannot identify an affirmative value.\claim{C06}

Evidence citations retain exact supporting source spans and declared context
dependencies. Temporal certainty is conditional on one selected state and its
episode assumptions. Deterministic selection among equally scored assignments
is not evidence that those assignments agree on the answer. Evaluation of
uncertainty across plausible assignments remains future work.

\section{Shared score function and inference comparisons}
\label{sec:inference}
Let $u_m(r)$ be a reading's supplied unary score minus its configured local
violation penalties. A non-null link has net score
$w(\ell)=\beta\,\mathrm{score}(\ell)-\mathrm{penalty}(\ell)$, while $n_m$ is
the supplied null score, without $\beta$ scaling. Over semantically feasible
assignments, the historical-interpretation objective is
\begin{equation}
J_{\mathrm{hist}}(r,\ell)=\sum_m u_m(r_m)
+\sum_{m:\ell_m\ne\varnothing}w(\ell_m)
+\sum_{m:\ell_m=\varnothing}n_m. \label{eq:historical}
\end{equation}
It scores interpretations of historical assertions, including assertions whose
support is later withdrawn. Scores are finite numeric inputs, not automatically
calibrated probabilities. All methods receive the same candidate graph, scores,
null alternatives, penalties, correction policy, and feasibility semantics.
\claim{C07}

The \emph{independent-reading baseline} first selects
$r_m=\arg\max_{r\in R_m}u_m(r)$ and then exactly optimizes all eligible link/null
assignments conditional on those readings. This is stronger than independently
choosing each link. \emph{Iterative inference} starts from the same readings,
tries each alternative at one mention, exactly reoptimizes all links for each
trial, and accepts strict objective improvements. The restarted variant repeats
this procedure from seeded random reading assignments. \emph{Exact joint
inference} enumerates reading assignments and feasible link assignments to
maximize the same objective. Tie-breaking is deterministic.\claim{C08}

The shared default ceiling is eight mentions, five readings per mention,
250,000 conservatively counted states, 50 coordinate sweeps, and ten optional
restarts. The state bound is
$\prod_m|R_m|\prod_m(1+|L_m|)$, which overcounts links incompatible with
particular readings. Every method rejects a graph above that bound. Exact
search certifies the implemented finite objective optimum, not factual truth;
these common limits do not measure scalability beyond the reference regime.

Coupling local interpretation scores with temporal constraints is established
structured prediction~\cite{han-etal-2019-joint,ning-etal-2018-joint}. A standard
constrained-prediction adaptation is therefore an appropriate future solver
baseline for this graph. The present enumerator supplies a small-instance
reference, not evidence of a new optimization algorithm.\claim{C14}

\subsection{An active-support ablation and its failure condition}
\label{sec:objective}
Let $A$ contain active resolved mentions, $U$ the selected unresolved mentions,
and $b_m=\max_{r\in R_m\text{ unresolved}}u_m(r)$. Let $T$ be selected temporal
links with both endpoints active, $I$ the remaining selected temporal links,
$C$ the operative correction links, and $N$ mentions choosing null. Write
$\operatorname{out}(\ell)$ for a link's outgoing mention. The
implemented ablation is
\begin{align}
J_{\mathrm{active}}={}&\sum_{m\in A\cup U}u_m(r_m)
+\sum_{m\notin A\cup U}b_m
+\sum_{\ell\in T\cup C}w(\ell) \nonumber\\
&+\sum_{\ell\in I}n_{\operatorname{out}(\ell)}+\sum_{m\in N}n_m.
\label{eq:active}
\end{align}
Inactive support is anchored to existing unresolved/null references, avoiding
dependence on arbitrary additive unary and outgoing net-score offsets. This
does not provide calibration, duplication invariance, or invariance to changing
only the unary or link scale.\claim{C09}

The ablation can suppress a valid correction. With fixed readings, if selecting
a correction $c\to q$ only withdraws target $q$, its objective gain over null is
\begin{equation}
\Delta J_{\mathrm{active}}=(w(c\to q)-n_c)-(u_q(r_q)-b_q).
\end{equation}
A positive historical reading margin can therefore outweigh correction support.
Conversely, withdrawing a negative-margin reading can increase the score.
Existing authored scores expose this failure in development replays, including
a source-record correction example. Scores were not retuned to remove the
failure. The historical objective remains the default; this ablation is not
presented as a successful method.\claim{C09}

\section{Data and score provenance: present evidence}
\label{sec:data}
The completed v0.6 natural pilot uses two Nike records concerning a CEO
transition, within one development history and one operational information
prefix. A source-only conversational extraction produced an invalid initial
temporal representation; one recorded repair separated actual observations
from planned information. The retained graph has seven anchored mentions,
seven resolved readings, seven unresolved alternatives, and eleven candidate
links. Each mention has only one resolved reading. This pilot therefore does
not test selection among competing resolved factual interpretations.\claim{C10}

A separate conversational scoring pass supplied 32 ordinal judgments:
resolved and unresolved readings, candidate links, and null choices. The backend
revision was unavailable, so these outputs are explicitly unversioned and are
not log probabilities. Exact requests, outputs, repair lineage, and score-cache
hashes are retained. Replaying stored outputs is reproducible artifact
processing; it is not reproducible generation from a pinned model.\claim{C10}

All four decoders selected the same readings and links under both objectives;
the selected states had no correction or essential-support inactivity. Three
declared monotone score recodings also produced unchanged selections in twelve
historical-objective decodes. No answer labels or factual accuracy measurement
were used in these runs. Agreement on this example and these recodings provides
neither a joint-inference gain nor general score-robustness evidence.\claim{C10}

Engineering diagnostics separately check schema, temporal closure, corrections,
shared-objective evaluation, and artifact replay. The v0.6 objective audit
reports 23,497 checks without a mismatch against its independent finite oracle.
These checks support implementation contracts within their enumerated domains;
they are not independent natural-data observations or a measure of extraction
quality.\claim{C11}

\section{Prospective evaluation and decision gates}
\label{sec:evaluation}
The next empirical unit is a prospectively sampled batch of twenty development
histories under the controlled document-stream protocol, with routine cases and
an enriched ambiguity/correction slice reported separately. Previously inspected
Nike, OpenAI, and Lyft examples remain
development demonstrations. A sampling rule, exclusion log, source version and
availability policy, delivery batches, and candidate budgets must be fixed before
extraction or method-output inspection. Simultaneous or uncertain ordering stays
in a common delivery batch. Every method receives identical source IDs and bytes
for a prefix. A present-day retrieval may contain later edits; controlled-stream
admission does not certify it as the original historical version. Existing
cutoffs and results are not retroactively admitted to the new benchmark.
Connected histories must remain together in subsequent
train, development, and held-out splits. This is a proposed study, not a
completed benchmark.\claim{C12}

Candidate generation should be frozen without questions, answer labels, or
decoder outputs. The evaluation must distinguish missing valid alternatives,
unsupported alternatives, structural exclusions, and scoring errors. A useful
controlled comparison can generate hypotheses from local windows while giving
both structured inference and a strong contextual/listwise baseline identical
full-prefix access. Full-context source judgments evaluate the predictions;
they must not act as an oracle that removes every falsifiable hypothesis before
inference. Residual semantic ambiguity should be reported separately from
extractor uncertainty. These principles are frozen in the prospective stream
protocol; the sampling frame, annotation materials, and executable collection
pipeline still need to be completed before the batch is run.\claim{C12}

Independent source-based reference judgments should cover interpretations,
update targets, temporal support, and answerability, with recorded provenance
and a disagreement process. Model-assisted judgments must be labeled as such;
isolated model roles do not create human gold. Candidate coverage, link validity,
factual answers, supporting evidence, indeterminate/abstain behavior, and cost
must be reported separately from objective gaps. Natural-history resampling and
uncertainty estimates should use independent histories rather than treating
correlated questions as independent samples.\claim{C12}

The comparison must include contextual extraction, independent-reading plus
exact conditional links, coordinate ascent, restarted inference, and the exact
reference where feasible. Model sizes, prompts, objectives, and score scales
are separate ablations. Record every model call, token budget, failure, wall
time, and measured memory. If nontrivial coupling survives source adjudication
and strong contextual/restarted baselines, a scalable algorithm should target
that measured structure. If it does not, the paper direction must narrow to a
supported representation or uncertainty question after a fresh novelty audit.
A positive outcome is not guaranteed.\claim{C12}

\section{Related work and the remaining empirical question}
\label{sec:relatedwork}
\paragraph{Temporal memory and evolving evidence.}
Source-grounded temporal graph memory, separated clocks, and typed updates have
close precedents. SodaMem represents fact events with evidence spans, modality,
status, and multiple temporal fields~\cite{wan2026sodamem}. Zep separates
transaction and fact-validity times while resolving entities and updating
contradictory facts~\cite{rasmussen2025zep}. Temporal Semantic Memory separates
semantic occurrence from dialogue time and maintains durative
facts~\cite{su2026temporal}. A-TMA preserves states and audits suspicious
same-slot pairs through proposal, judgment, and commit
stages~\cite{shi2026atma}. These overlaps make provenance, multiple clocks,
and state-preserving ingestion insufficient novelty claims. The listed arXiv
works are cited as preprints; no external reported result is reproduced here.

TG-RAG studies temporal graph updates and evaluation over an incrementally
enlarged corpus~\cite{han2025tgrag}. HoH evaluates evolving factual questions
and interference from obsolete evidence using Wikipedia
versions~\cite{ouyang-etal-2025-hoh}. A controlled document stream or a benchmark
of changing facts is therefore not a new contribution by itself. Our proposed
evaluation must make its additional distinctions explicit: correcting an
erroneous assertion versus changing a previously valid state, announced versus
observed outcomes, and uncertainty about interpretations and update targets.
Whether those distinctions yield useful new evidence remains untested.

\paragraph{Joint constrained prediction.}
Han et al. jointly predict event labels and temporal relations using shared
representations and constrained structured prediction, allowing relation
evidence to revise an event decision~\cite{han-etal-2019-joint}. Ning et al.
couple temporal and causal relations through joint constraints and integer
linear programming~\cite{ning-etal-2018-joint}. These are direct precedents for
jointly choosing local labels and relations under an additive score. Their
different application does not establish an expressive limitation. Our proposed
source-claim alternatives, correction-dependent activation, and uncertain
intervals require a formal comparison and a measured benefit over a
straightforward constrained-prediction adaptation.

The narrow working question is whether correction-dependent reconstruction
benefits from retaining and revisiting competing interpretations once their
support and temporal consequences interact. A fair experiment must compare
strong full-context extraction, source-based verification, temporal-memory
update auditing, and standard constrained inference with matched evidence.
If correction semantics help while all solvers select the same readings, the
finding concerns representation or update policy rather than joint search.
Neither a new task nor a new algorithm is established by this draft.\claim{C14}

\section{Limitations and literature work still required}
\label{sec:limitations}
The implementation uses connected intervals at day resolution, claim-level
modality, single-valued qualified keys, one primary outgoing link per mention,
and small exhaustive components. It lacks endpoint-specific modality,
alternative-assignment uncertainty, broad correction actions, authenticated
historical availability, and a scalable solver. Exact quotes and hashes verify
integrity rather than entailment. Pretraining exposure and undeclared context
cannot be ruled out by an isolated request. Existing development data do not
establish natural candidate recall, factual accuracy, or generalization.
\claim{C13}

The targeted review above is not exhaustive. Remaining literature gaps include
joint entity resolution and event coreference, modern temporal relation
extraction, probabilistic temporal databases, assumption-based truth maintenance,
belief revision, temporal knowledge-base population, and uncertainty-aware
structured prediction. Later versions and released code of the cited systems
also require inspection before a formal novelty claim. The comparison must
align task assumptions, supervision, evidence access, and evaluation rather
than infer a gap from different terminology or omitted search results.
\todo{verified primary sources and closest-method comparison for these remaining gaps}
Classical closure, exhaustive search, withdrawal semantics, and reproducible
caching alone do not establish an algorithmic contribution. The intended ACL
2027 target is a research objective, not evidence of novelty or acceptance.
\claim{C14}

\bibliographystyle{plain}
\bibliography{paper/references_v07}

\end{document}
